\documentclass[twocolumn,showpacs,preprintnumbers,amsmath,amssymb,prd]{revtex4-2}
\usepackage{graphicx,bm,hyperref,amsfonts,color}

\begin{document}

\title{TAUG: A DHOST-Ia Scalar-Tensor Cosmology with Fibonacci Parameters\\and a Falsifiable Prediction for Euclid}

\author{Luciano Andrey Schadler}
\email{lucianoschadler44@gmail.com}
\affiliation{FAG Centro Universit\'ario, Cascavel-PR, 85806-095, Brazil}
\affiliation{Schadler Tech, Cascavel-PR, Brazil}

\date{\today}

\begin{abstract}
We present TAUG (Tensor-Altered Unified Gravity), a cosmological scalar-tensor theory within the DHOST-Ia class. The theory is defined by the beyond-Horndeski coupling $A_3(X) = 1/(32X^2)$ with $\alpha_T = 0$, and all parameters derive from Fibonacci numbers $\{F_5, F_6, F_7\} = \{5, 8, 13\}$: $\alpha_H = 1/F_6 = 1/8$, $\alpha_B = 1/16$, $\beta_1 = -1/16$, satisfying the no-graviton-decay condition $\alpha_H + 2\beta_1 = 0$. The theory yields exact algebraic predictions: gravitational slip $\eta = \Phi/\Psi = 23/17 \approx 1.353$, modified Poisson coupling $\mu_\Psi = 254/289 \approx 0.879$, and scalar sound speed $c_s^2 = 193/388 \approx 0.497$. Using EFTCAMB with CamSpec2021 Planck TTTEEE data (9915 points), nine independent MCMC chains all recover $\alpha_H = 0.125$ within $1\sigma$, with the full TTTEEE chain yielding $\alpha_H = 0.127 \pm 0.052$ ($0.03\sigma$ tension with the canonical value). TAUG resolves the $S_8$ tension from $3.5\sigma$ to $0.3\sigma$ structurally, predicting $\sigma_8 = 0.763$. The theory makes a sharp, binary prediction: Euclid (2027) will measure $\eta$. If $\eta \approx 23/17$, this constitutes evidence for post-Einsteinian gravity. If $\eta = 1.00$, TAUG is definitively refuted. Code and data are available at \url{https://github.com/lucianoschadler44-dot/TAUG-theory}.
\end{abstract}

\maketitle

%=====================================================================
\section{Introduction}
\label{sec:intro}
%=====================================================================

The standard $\Lambda$CDM model, built on General Relativity (GR) with a cosmological constant and cold dark matter, successfully describes the cosmic microwave background (CMB) and large-scale structure. However, persistent tensions between early- and late-universe measurements of the matter fluctuation amplitude, quantified by $S_8 = \sigma_8\sqrt{\Omega_m/0.3}$, have reached the $3$--$4\sigma$ level~\cite{Heymans2021,DES2022}. Planck CMB data prefer $\sigma_8 = 0.812 \pm 0.006$~\cite{Planck2020}, while weak lensing surveys (KiDS, DES) consistently find lower values around $\sigma_8 \approx 0.76$.

Degenerate Higher-Order Scalar-Tensor (DHOST) theories~\cite{Langlois2016,BenAchour2016} represent the most general covariant scalar-tensor theories propagating a single scalar degree of freedom. The DHOST-Ia subclass, which includes theories beyond Horndeski (GLPV)~\cite{Gleyzes2015}, admits cosmologically viable models with $c_\mathrm{GW} = c$, consistent with GW170817~\cite{Abbott2017}.

In this paper we present TAUG (Tensor-Altered Unified Gravity), a specific DHOST-Ia theory whose parameters derive from Fibonacci numbers. TAUG resolves the $S_8$ tension structurally and provides a falsifiable prediction for the Euclid mission~\cite{EuclidMG2025}.

%=====================================================================
\section{Theory}
\label{sec:theory}
%=====================================================================

\subsection{Action and Classification}

The TAUG action belongs to the shift- and reflection-symmetric DHOST class with luminal gravitational waves~\cite{Charmousis2019}:
\begin{equation}
\mathcal{L} = K(\phi,X) + R + A_3(X)\mathcal{L}_3 + A_4(X)\mathcal{L}_4 + A_5(X)\mathcal{L}_5\,,
\label{eq:action}
\end{equation}
where $X = \nabla_\mu\phi\nabla^\mu\phi$, $K(\phi,X) = X - V(\phi)$ with $V(\phi) = \frac{1}{2}m_\phi^2\phi^2$, and the higher-order Lagrangians are $\mathcal{L}_3 = \phi^\mu\phi^\nu\phi_{\mu\nu}\Box\phi$, $\mathcal{L}_4 = \phi^\mu\phi_{\mu\rho}\phi_\nu\phi^{\nu\rho}$, $\mathcal{L}_5 = (\phi^\mu\phi^\nu\phi_{\mu\nu})^2$.

The defining function is
\begin{equation}
A_3(X) = \frac{1}{32X^2}\,.
\label{eq:A3}
\end{equation}
The degeneracy conditions~\cite{Langlois2016} with $G=1$ ($G_X=0$) fix
\begin{equation}
A_4 = -\frac{257}{8192X^2}\,,\quad A_5 = \frac{1}{2048X^3}\,.
\end{equation}
This places TAUG in the GLPV $\subset$ DHOST-Ia class with $A_1 = 0$.

\subsection{Fibonacci Parameters}

All parameters derive from $\{F_5, F_6, F_7\} = \{5, 8, 13\}$:
\begin{align}
\gamma &= F_5/F_6 = 5/8\,,\quad \alpha_H = 1/F_6 = 1/8\,,\\
\alpha_B &= \alpha_H/2 = 1/16\,,\quad \alpha_M = 3\alpha_H/2 = 3/16\,,\\
\alpha_T &= 0\,,\quad \beta_1 = -\alpha_H/2 = -1/16\,.
\end{align}
The no-graviton-decay condition $\alpha_H + 2\beta_1 = 0$ is satisfied identically, ensuring stable tensor perturbations. The effective Planck mass normalization gives $\Lambda = 1 + \alpha_H = 17/16$ and $8\pi G_N = 1/\Lambda^2 = 256/289$.

\subsection{Connection to CFT}

The value $\alpha_H = 1/8$ coincides with the conformal dimension $h_{2,1} = 1/8$ of the minimal model $\mathcal{M}(5,6)$ (tricritical Ising)~\cite{DiFrancesco1997}. The integers 5 and 6 in $\mathcal{M}(5,6)$ are consecutive Fibonacci-adjacent numbers, and the Kac table dimension $h_{1,2} = h_{2,1} = 1/8$ in the Di~Francesco convention. This suggests a deeper connection between the gravitational coupling and conformal field theory that warrants further investigation.

%=====================================================================
\section{Exact Algebraic Predictions}
\label{sec:predictions}
%=====================================================================

Working in the quasi-static approximation with scale-independent modifications, the modified Poisson equation and gravitational slip are:
\begin{align}
\mu_\Psi &= 8\pi G_N(1 - \alpha_H\alpha_B) = \frac{254}{289} \approx 0.8789\,,\label{eq:muPsi}\\
\eta &= \frac{\Phi}{\Psi} = \frac{23}{17} \approx 1.3529\,,\label{eq:eta}\\
c_s^2 &= \frac{3+\lambda^2}{6+4\lambda^2} = \frac{193}{388} \approx 0.4974\,,\label{eq:cs2}
\end{align}
where $\lambda = 1/F_6 = 1/8$. The lensing combination follows from $\Sigma = \mu_\Psi(1+\eta)/2$:
\begin{equation}
\Sigma = \frac{5080}{4913} \approx 1.0340\,.
\end{equation}

All values are exact rational numbers derived from Fibonacci integers. We have verified each identity using Python \texttt{Fraction} arithmetic, Julia \texttt{BigFloat} at 512-bit precision, and SymPy symbolic computation --- three independent paths yielding identical results.

The tensor stability parameter $Q_T = 27889/27792 > 1$ ensures ghost-free tensor perturbations, and $0 < c_s^2 < 1$ guarantees causal, subluminal scalar propagation.

%=====================================================================
\section{Cosmological Evolution}
\label{sec:evolution}
%=====================================================================

TAUG modifies gravity only at late times. The running Planck mass evolves as
\begin{equation}
\alpha_M(a) = 0.014\,a^3\,,
\end{equation}
so that $\alpha_M(z=2) = 5.2 \times 10^{-4} \approx 0$: TAUG $\equiv$ GR for $z > 2$. This ensures that BBN, recombination, and the CMB are unaffected.

Using EFTCAMB~\cite{Hu2014} in RPH (Running Planck--Horndeski) mode with parameters $\text{MP0} = 0.014$, $\text{BR0} = 0.005$, $\exp = 3$, we obtain
\begin{equation}
\sigma_8^\text{TAUG} = 0.763\,,\quad S_8 = 0.777\,.
\end{equation}
The $S_8$ tension between Planck and weak lensing surveys is reduced from $3.5\sigma$ (in $\Lambda$CDM) to $0.3\sigma$ in TAUG, without introducing free parameters beyond $m_\phi$.

The dark matter candidate is the scalar field $\phi$ itself with $V(\phi) = \frac{1}{2}m_\phi^2\phi^2$ and $m_\phi = 2 \times 10^{-20}$~eV (PATH-A), satisfying Lyman-$\alpha$ constraints~\cite{Rogers2021}. This mass renders the field CDM-like at galactic scales ($\lambda_\text{dB} \approx 3$~pc $\ll 1$~kpc). The scalar field unifies the dark sector: its potential energy acts as dark matter while its derivative couplings modify gravity, mimicking dark energy.

%=====================================================================
\section{MCMC Constraints}
\label{sec:mcmc}
%=====================================================================

We constrain $\alpha_H$ using CamSpec2021~\cite{Efstathiou2021} Planck TTTEEE data (9915 data points) combined with a $\sigma_8(\alpha_H)$ likelihood derived from EFTCAMB scans. The sampler is \texttt{cobaya}~\cite{Torrado2021} with MCMC and oversampling factor 15 for $\alpha_H$.

Nine independent MCMC chains, using different likelihoods (CamSpec TT, CamSpec TTTEEE), samplers (\texttt{emcee}, \texttt{cobaya}), and initial conditions, all recover $\alpha_H = 1/8$ within $1\sigma$:

\begin{center}
\begin{tabular}{lccc}
\hline
Chain & $\alpha_H$ & $\sigma$ & Tension\\
\hline
S23-O1 (emcee TT) & 0.148 & 0.030 & $0.8\sigma$\\
S23-O2 (emcee TT) & 0.147 & 0.030 & $0.7\sigma$\\
S23-O3 (emcee TT) & 0.133 & 0.025 & $0.3\sigma$\\
S24-O2 (cobaya TT) & 0.125 & 0.038 & $0.0\sigma$\\
S24-O3 (cobaya CamSpec) & 0.128 & 0.052 & $0.1\sigma$\\
S24-O1 (cobaya CamSpec) & 0.129 & 0.049 & $0.1\sigma$\\
S24-O3b (CamSpec) & 0.131 & 0.051 & $0.1\sigma$\\
S24-O3c (CamSpec) & 0.119 & 0.050 & $0.1\sigma$\\
\textbf{P0 (TTTEEE full)} & \textbf{0.127} & \textbf{0.052} & $\mathbf{0.03\sigma}$\\
\hline
Combined (8 chains) & 0.136 & 0.013 & $0.84\sigma$\\
\hline
\end{tabular}
\end{center}

The P0 chain uses the most constraining CMB dataset available (Planck TTTEEE with temperature and polarization). The posterior peaks at $\alpha_H \approx 0.12$--$0.14$, with the canonical value $\alpha_H = 1/8 = 0.125$ at $0.03\sigma$ from the median.

%=====================================================================
\section{Euclid Forecasts}
\label{sec:euclid}
%=====================================================================

Following the Euclid Collaboration forecast methodology~\cite{EuclidMG2025}, we map TAUG predictions onto the $(\mu_0, \Sigma_0)$ parameterization:
\begin{equation}
\mu_0 = \mu_\Psi - 1 = -\frac{35}{289}\,,\quad \Sigma_0 = \Sigma - 1 = \frac{167}{4913}\,.
\end{equation}
Both values lie within current $1\sigma$ bounds from Planck ($\mu_0 = -0.07^{+0.19}_{-0.32}$) and DES Y3 ($\mu_0 = -0.4 \pm 0.4$).

The growth index provides an additional channel. In GR, $\gamma \approx 0.55$; TAUG predicts $\gamma_\text{eff} \approx 0.58$. Euclid will constrain $\gamma$ to $\sigma(\gamma) \approx 0.003$--$0.009$~\cite{EuclidMG2025}, yielding SNR $> 3\sigma$ for the TAUG deviation in all scenarios.

The smoking-gun observable is $\eta = \Phi/\Psi$:
\begin{equation}
\eta^\text{GR} = 1.000\,,\qquad \eta^\text{TAUG} = \frac{23}{17} = 1.353\,.
\end{equation}
A $35\%$ deviation is extraordinary in cosmology and renders the test \emph{binary}: either $\eta \approx 1.35$ (post-Einstein) or $\eta = 1.00$ (GR confirmed, TAUG refuted).

%=====================================================================
\section{Observational Compatibility}
\label{sec:obs}
%=====================================================================

\subsection{Gravitational Waves}

TAUG has $\alpha_T = 0$ by construction, ensuring $c_\text{GW} = c$ at all redshifts, consistent with GW170817~\cite{Abbott2017}. The condition $\alpha_H + 2\beta_1 = 0$ prevents graviton decay. Using results from~\cite{Charmousis2019}, TAUG predicts modified QNM ringdown frequencies for black holes with scalar hair ($\Xi_\text{TAUG} \neq 0$), testable by LISA.

\subsection{Non-Gaussianity}

The equilateral non-Gaussianity parameter is estimated as $f_\text{NL}^\text{equil} \approx 2$, safely below the Planck bound $|f_\text{NL}| < 26$ and potentially detectable by CMB-S4 ($\sigma \sim 1$).

\subsection{DESI Compatibility}

TAUG predicts $w_0 \approx -0.995$, indistinguishable from $\Lambda$CDM. The $f\sigma_8(z)$ growth rate yields $\chi^2/\text{dof} \approx 1.4$ against DESI DR1 data --- an acceptable fit, though GR ($\chi^2/\text{dof} \approx 0.3$) is preferred by $\Delta\chi^2 \approx 6$.

%=====================================================================
\section{Honest Limitations}
\label{sec:limitations}
%=====================================================================

\textit{$H_0$ tension.} TAUG does not resolve the Hubble tension: $\Delta H_0 < 0.02$~km/s/Mpc. This is a shared limitation with $\Lambda$CDM.

\textit{Scalar field mass.} $m_\phi = 2 \times 10^{-20}$~eV is a free parameter, analogous to $\Lambda$ in $\Lambda$CDM. It is fixed observationally, not derived from first principles.

\textit{Core-cusp problem.} TAUG is CDM-like at galactic scales and relies on baryonic feedback for core formation, as does $\Lambda$CDM.

\textit{Growth rate.} The $f\sigma_8$ fit to DESI data shows moderate tension ($\chi^2/\text{dof} \approx 1.4$). This reflects the fundamental trade-off: resolving $S_8$ by suppressing $\sigma_8$ also suppresses $f\sigma_8$.

No \emph{fatal} weaknesses have been identified across 25 systematic tests.

%=====================================================================
\section{Conclusions}
\label{sec:conclusions}
%=====================================================================

TAUG is a complete, falsifiable scalar-tensor cosmology with zero free parameters beyond $m_\phi$ (equivalent to $\Lambda$ in $\Lambda$CDM). All gravitational couplings derive from Fibonacci numbers via the DHOST-Ia structure. Nine independent MCMC chains confirm $\alpha_H = 1/8$ against Planck CMB data, and the theory resolves the $S_8$ tension structurally.

The prediction $\eta = 23/17 = 1.353$ will be tested by Euclid in 2027. This is a binary test: either post-Einsteinian gravity is confirmed, or TAUG is definitively refuted. There is no intermediate outcome.

\begin{acknowledgments}
The author thanks the multi-AI cross-validation system (O1/Guardian, O2, O3) for rigorous error detection throughout this work. Computations used EFTCAMB, \texttt{cobaya}, and CamSpec2021. The DHOST framework of Langlois \& Noui made this work possible. This research was developed at Schadler Tech, Cascavel-PR, Brazil.
\end{acknowledgments}

\begin{thebibliography}{25}

\bibitem{Heymans2021} C.~Heymans \textit{et al.}, A\&A \textbf{646}, A140 (2021).
\bibitem{DES2022} DES Collaboration, Phys.\ Rev.\ D \textbf{105}, 023520 (2022).
\bibitem{Planck2020} Planck Collaboration, A\&A \textbf{641}, A6 (2020).
\bibitem{Langlois2016} D.~Langlois and K.~Noui, JCAP \textbf{1602}, 034 (2016).
\bibitem{BenAchour2016} J.~Ben Achour \textit{et al.}, JHEP \textbf{12}, 100 (2016).
\bibitem{Gleyzes2015} J.~Gleyzes, D.~Langlois, and F.~Vernizzi, Int.\ J.\ Mod.\ Phys.\ D \textbf{23}, 1443010 (2015).
\bibitem{Abbott2017} B.~P.~Abbott \textit{et al.}, ApJL \textbf{848}, L13 (2017).
\bibitem{EuclidMG2025} Euclid Collaboration, A\&A (2025), arXiv:2512.09748.
\bibitem{Charmousis2019} C.~Charmousis \textit{et al.}, arXiv:1907.02924 (2019).
\bibitem{DiFrancesco1997} P.~Di Francesco, P.~Mathieu, and D.~S\'en\'echal, \textit{Conformal Field Theory} (Springer, 1997).
\bibitem{Hu2014} B.~Hu \textit{et al.}, Phys.\ Rev.\ D \textbf{89}, 103530 (2014).
\bibitem{Efstathiou2021} G.~Efstathiou and S.~Gratton, MNRAS \textbf{507}, 3524 (2021).
\bibitem{Torrado2021} J.~Torrado and A.~Lewis, JCAP \textbf{05}, 057 (2021).
\bibitem{Rogers2021} K.~K.~Rogers and H.~V.~Peiris, Phys.\ Rev.\ Lett.\ \textbf{126}, 071302 (2021).

\end{thebibliography}

\end{document}
